\documentclass{article}

\usepackage[T1]{fontenc}
\usepackage[utf8]{inputenc}
\usepackage{hyperref}
\usepackage{url}
\usepackage{booktabs}
\usepackage{nicefrac}
\usepackage{microtype}
\usepackage{xcolor}

%%%

\usepackage{subcaption}
\usepackage{graphicx}
\usepackage{multirow}
\usepackage{amsmath,amssymb,amsfonts}
\usepackage{amsthm}
\usepackage{mathrsfs}
\usepackage{xcolor}
\usepackage{textcomp}
\usepackage{manyfoot}
\usepackage{algorithm}
\usepackage{algorithmicx}
\usepackage{algpseudocode}
\usepackage{listings}

\newtheorem{theorem}{Theorem} % continuous numbers
%%\newtheorem{theorem}{Theorem}[section] % sectionwise numbers
%% optional argument [theorem] produces theorem numbering sequence instead of independent numbers for Proposition
\newtheorem{proposition}[theorem]{Proposition}% 
\newtheorem{lemma}{Lemma}% 
%%\newtheorem{proposition}{Proposition} % to get separate numbers for theorem and proposition etc.

\newtheorem{example}{Example}
\newtheorem{remark}{Remark}

\newtheorem{definition}{Definition}
\newtheorem{assumption}{Assumption}

%%%


\title{Calibrated Perplexity for Robust AI-Generated Text Detection Across Languages and Genres}


% The \author macro works with any number of authors. There are two commands
% used to separate the names and addresses of multiple authors: \And and \AND.
%
% Using \And between authors leaves it to LaTeX to determine where to break the
% lines. Using \AND forces a line break at that point. So, if LaTeX puts 3 of 4
% authors names on the first line, and the last on the second line, try using
% \AND instead of \And before the third author name.


\author{
  Vladimir Chaikin\\
  MIPT\\
  Moscow, Russia\\
  \texttt{email@phystech.edu}\\
  \And
  Anastasia Voznyuk\\
  MIPT\\
  Moscow, Russia\\
  \texttt{email@phystech.edu}\\
}


\begin{document}


\maketitle

\begin{abstract}
    This paper invistigates the problem of robust AI-generated text detection under distribution shift. While simple perplexity is a simple solution for such detection, its oppurtunities lose significantly on out-of-domain data. We propose a novel method - calibrated perplexity - which normalizes raw perplexity scores using domain-specific language model baselines. This approach constructs a stable feature space less sensitive to text length, genre, and language variations. Experiments across multiple domains and languages are going to demonstrate improved in-domain and out-of-domain detection accuracy compared to standard perplexity-based methods.
\end{abstract}

\textbf{Keywords:} Calibrated perplexity
Cross-domain detection
Robustness
Machine-generated text
Text classification
Multilingual detection
Genre adaptation
Statistical methods
Language modeling

\section{Introduction}\label{sec:intro}

Fast improvement and widespread accessibility of Large Language Models (LLMs) have destroyed the lines between human-written and machine-generated text. While this technology offers benefits, it also presents significant disadvantages, including academic dishonesty, disinformation, and the fall of trust in digital content \cite{xu2024detecting, hans2024spotting}. This has created an inevitable need for robust methods to detect AI-generated content.

Current detectors often rely on fine-tuned classifiers or statistical signatures taken from the text \cite{gehrmann2019gltr}. Useful and computationally efficient feature is perplexity, a measure of how "surprising" a text is to a language model. The assumption is that LLMs, which are trained to maximize likelihood, will generate text with lower perplexity compared to human writing. As it is shown in \cite{bakhteev2022automatic}, a significant change comes from distribution shift. The perplexity of a text is not an absolute value; it is highly dependent on the domain (e.g., news, scientific articles, social media) and language it was trained on. A model fine-tuned for one genre often fails on another, causing unreliable detection in real world, where the target domain is unknown.

Recent works have attempted to solve this by developing zero-shot detectors that are more robust to unseen prompts or models. For instance, Binoculars \cite{hans2024spotting} uses a ratio of scores from two different LLMs for calibration. Similarly, research on multilingual and cross-domain detection highlights the necessity of adapting methods to diverse data distributions \cite{gritsay2023artificially}. Despite these advances, a method for calibrating perplexity for each domain itself to make it robust across different genres and languages remains still not-explored.

In this paper, we face the challenge of domain shift and try to avoid its effect on detection. We propose a novel method called calibrated perplexity, which normalizes perplexity scores against domain-specific baselines. The final goal is not to find out one universal calibrated perplexity score, but to get a set of parameters to have an opportunity to learn fastly new zero-shot models. Our approach includeses training KenLM n-gram models on text from multiple domains and using these models to compute a relative, rather than absolute, measure of text predictability. The core novelty of our work lies in transforming a domain-dependent score into a stable feature for classification. To check our method, we conduct experiments across multiple domains and languages, comparing in-domain and out-of-domain detection.


\section{Problem statement}

Let $\mathbf{W}$ be our alphabet, tokens consist of characters from that alphabet.

Let 
\[
\mathbb{D} = \left\{ \left[ t_j \right]_{j=1}^n \;\Big|\; t_j \in \mathbf{W},\ n \in \mathbb{N} \right\}
\]
be the space of our documents.

We have a set of $N$ documents
\[
\mathbf{D} = \bigcup_{i=1}^{N} D^i,\quad D^i \in \mathbb{D}.
\]

We assume that each document may contain fragments generated by one of $K$ text-generative models, or be fully written by a human. Accordingly, we introduce $K+1$ labels for multiclass classification, where $0$ denotes a human-written fragment and $\{1,\dots,K\}$ denote fragments generated by the corresponding language model. Let $\mathbf{C}$ be the set of these labels.

To capture the statistical properties of text fragments, we consider a probabilistic model of the language. For a given fragment $F = [t_a, t_{a+1}, \dots, t_b]$ we define its probability as the product of conditional probabilities of each token given its $n-1$ predecessors:
\[
P(F) = \prod_{i=a}^{b} P(t_i \mid t_{i-1}, \dots, t_{i-n+1}).
\]

The choice of such an $n$-gram model is motivated by Hypothesis 1 in \cite{gritsay2023artificially}, which states that generative language models tend to sample from the head of the token distribution, i.e., they prefer high-probability continuations. This makes $n$-gram statistics a natural baseline for distinguishing human-written from machine-generated text.

The log-probability of a fragment is then
\[
\log P(F) = \sum_{i=a}^{b} \log P(t_i \mid t_{i-1}, \dots, t_{i-n+1}),
\]
and the perplexity is defined as the inverse geometric mean of the token probabilities:
\[
\text{PPL}(F) = P(F)^{-\frac{1}{|F|}} = 2^{-\frac{1}{|F|}\log_2 P(F)},
\]
where $|F| = b-a+1$ is the length of the fragment in tokens.

Perplexity quantifies how "surprising" a fragment is under the language model: lower values indicate that the fragment is more predictable.

Given a training set of documents $\mathbf{D}_{\text{train}} \subset \mathbb{D}$ with known labels, we learn the conditional probabilities $P(t_i \mid t_{i-1}, \dots, t_{i-n+1})$ by maximizing the likelihood of the training data. This leads to the following optimization problem:
\[
\mathcal{L}(\boldsymbol{\theta}) = -\sum_{F \in \mathbf{D}_{\text{train}}} \log P(F) \rightarrow \min_{\boldsymbol{\theta}},
\]
where $\boldsymbol{\theta}$ denotes the parameters of the $n$-gram model (the conditional probability tables).

The training objective is the negative log-likelihood of the observed documents, which is equivalent to minimizing the cross-entropy between the model distribution and the empirical distribution of the training data. This optimization is typically performed using count-based estimation with appropriate smoothing to handle sparse $n$-gram occurrences.

For classification, we introduce a threshold $\tau_d$ for each domain $d$. A fragment $F$ is classified as machine-generated if:
\[
\text{PPL}_{\mathcal{M}_d}(F) < \tau_d,
\]
and as human-written otherwise. The threshold $\tau_d$ is selected on a validation set to maximize the F1-score, balancing precision and recall.



\section{Method}\label{sec:method}

\subsection{Overview}

We propose an ensemble-based method for AI-generated text detection that operates without knowledge of the test text's domain. The method consists of three stages:

\begin{enumerate}
    \item \textbf{Training:} Train separate $n$-gram language models on different domains using only human-written texts.
    \item \textbf{Threshold calibration:} For each domain model, select a classification threshold that maximizes F1-score on a validation set from the same domain.
    \item \textbf{Ensemble inference:} For a new text, compute perplexities from all domain models and combine their predictions using weighted voting.
\end{enumerate}

\subsection{Notation}

Let $\mathcal{D} = \{d_1, d_2, \dots, d_M\}$ be a set of $M$ domains (e.g., news, essays, scientific texts). For each domain $d$, we have:
\begin{itemize}
    \item Training set $\mathbf{T}_d \subset \mathbb{D}$ consisting of human-written texts from domain $d$
    \item Validation set $\mathbf{V}_d$ with balanced human and machine texts from domain $d$
    \item Test set $\mathbf{X}_d$ held-out from the same distribution as validation
\end{itemize}

For a text $x$, let $\text{PPL}_{\mathcal{M}_d}(x)$ denote the perplexity computed by the language model trained on domain $d$.

\subsection{Model Training}

For each domain $d \in \mathcal{D}$, we train a trigram language model $\mathcal{M}_d$ using KenLM with modified Kneser-Ney smoothing. The training data consists exclusively of human-written texts from that domain:
\[
\mathcal{M}_d = \text{KenLM}(\mathbf{T}_d), \quad \text{where } \mathbf{T}_d = \{x \in \mathbb{D} \mid \text{source}(x) = \text{human}, \text{domain}(x) = d\}.
\]

The training objective is to minimize the negative log-likelihood of the training corpus:
\[
\mathcal{L}(\mathcal{M}_d) = -\sum_{x \in \mathbf{T}_d} \log P_{\mathcal{M}_d}(x) \rightarrow \min.
\]

\subsection{Threshold Calibration}

For each model $\mathcal{M}_d$, we find an optimal threshold $\tau_d$ that maximizes the F1-score on the validation set $\mathbf{V}_d$:
\[
\tau_d = \arg\max_{\tau} \text{F1}\left( \mathbf{1}[\text{PPL}_{\mathcal{M}_d}(x) < \tau],\; y(x) \right),
\]
where $y(x) = 1$ for machine-generated texts and $y(x) = 0$ for human-written texts.

The F1-score is defined as the harmonic mean of precision and recall:
\[
\text{F1} = 2 \cdot \frac{\text{precision} \cdot \text{recall}}{\text{precision} + \text{recall}}.
\]

In practice, we search over 200 evenly spaced thresholds between the minimum and maximum perplexity values observed on the validation set.

\subsection{Ensemble Inference}

For a test text $x$ with unknown domain, we compute perplexities $p_d = \text{PPL}_{\mathcal{M}_d}(x)$ for all $d \in \mathcal{D}$. These perplexities are then used in two complementary ensemble strategies.

\subsubsection{Strategy A: Machine-Oriented Ensemble (Unweighted Voting)}

This ensemble focuses on detecting machine-generated texts. Each model makes an individual prediction $h_d^{\text{machine}}(x) = \mathbf{1}[p_d < \tau_d]$, where $1$ indicates "machine-generated". The final decision is based on simple majority voting:
\[
H_{\text{machine}}(x) = \mathbf{1}\left[ \sum_{d \in \mathcal{D}} h_d^{\text{machine}}(x) \geq \lceil M/2 \rceil \right],
\]
where $M = |\mathcal{D}|$ is the number of domains. This strategy is effective when the input text is likely to be machine-generated, as it requires a simple majority of models to agree.

\subsubsection{Strategy B: Human-Oriented Ensemble (Perplexity-Weighted Voting)}

This ensemble focuses on detecting human-written texts. The key insight is that a text originating from a particular domain will yield minimal perplexity for the model trained on that domain. Therefore, we sort the domains by their perplexity in ascending order and assign decreasing weights:
\[
w_{(1)} \geq w_{(2)} \geq \dots \geq w_{(M)},\quad \sum_{k=1}^{M} w_k = 1,
\]
where $(k)$ denotes the index of the model with the $k$-th smallest perplexity. In our implementation, we use:
\[
[w_1, w_2, w_3, w_4, w_5, w_6, w_7] = [0.50, 0.20, 0.10, 0.07, 0.05, 0.04, 0.04].
\]

Each model makes an individual prediction $h_d^{\text{human}}(x) = \mathbf{1}[p_d \geq \tau_d]$, where $1$ indicates "human-written". The weighted sum is:
\[
S_{\text{human}}(x) = \sum_{k=1}^{M} w_k \cdot h_{(k)}^{\text{human}}(x),
\]
and the final decision is:
\[
H_{\text{human}}(x) = \mathbf{1}\left[ S_{\text{human}}(x) \geq \theta \right],
\]
where $\theta$ is a global threshold (typically $\theta = 0.5$). This strategy prioritizes models that best fit the input text (lowest perplexity), making it particularly sensitive to texts that are likely to be human-written.

\subsubsection{Combined Approach}

The two strategies complement each other:
\begin{itemize}
    \item Strategy A (unweighted majority) excels at catching machine-generated texts, as any sign of machine generation from multiple models is sufficient.
    \item Strategy B (perplexity-weighted) excels at confirming human-written texts, as a human text from a known domain will be recognized by its matching model.
\end{itemize}

In practice, we can either:
\begin{enumerate}
    \item Apply both strategies and combine their outputs (e.g., by averaging),
    \item Select the appropriate strategy based on the application requirements (e.g., use Strategy A for spam filtering, Strategy B for academic integrity checks),
    \item Use the weighted sum $S_{\text{human}}(x)$ directly as a confidence score.
\end{enumerate}

The optimal threshold $\theta$ for Strategy B is selected on the validation set by maximizing the balanced accuracy:
\[
\theta^* = \arg\max_{\theta} \text{BalancedAcc}\left( \mathbf{1}[S_{\text{human}}(x) \geq \theta],\; y(x) \right),
\]
where $y(x) = 1$ for human-written texts.


\subsection{Properties}

The proposed method has the following properties:

\begin{itemize}
    \item \textbf{Domain-agnostic inference:} No prior knowledge of the test text's domain is required.
    \item \textbf{Interpretability:} Each model provides a clear per-domain prediction; low perplexity indicates that the text "belongs" to that domain.
    \item \textbf{Robustness:} Weighted voting reduces the impact of models that perform poorly on the input text.
    \item \textbf{Scalability:} Adding a new domain requires training one additional KenLM model and computing its validation accuracy.
\end{itemize}

\subsection{Optimization Problem}

\[
\sum_{i=1}^{N_d}\sum_{x_j \in X_i}(M_{\mu_i,i}(x) - y_{ij})^2 + \sum_{i < j\leq N_d}(\text{TH}_{M_{\mu_i,i}}(x) - \text{TH}_{M_{\mu_j,j}}(x))^2 \to \min_{\bar{\mu}}
\]
\small
We minimize the variance of thresholds in ensemble across different domains for the same text $x$, making the score domain-invariant

Our ensemble approach approximates this ideal by weighting models according to their reliability. When all models give similar perplexities, the confidence is high; when they disagree, the weighted sum provides a robust decision.

\section{Experiments}\label{sec:exp}

\subsection{Experimental Setup}

We conduct experiments on a collection of Russian-language texts spanning seven domains: news, scientific texts, stories, articles, and 3 "hard" domains (hard\_news, hard\_science, hard\_essays). The dataset is balanced with approximately 3,500 human-written and 3,500 machine-generated texts per domain. Texts shorter than 50 words are filtered out. All models are trained using KenLM with trigram order ($n=3$) and modified Kneser-Ney smoothing. 

\subsection{Experiment 1: Single-Domain Training and In-Domain Testing}

For each domain, we train a KenLM model exclusively on human-written texts from that domain. We then evaluate the model on a held-out test set from the same domain, using a threshold optimized by F1-score on a validation set.

\textbf{Results:} In-domain performance varies significantly across domains, with F1-scores ranging from 0.85 (AINL science) to 0.53 (reviews). The highest balanced accuracy (0.85) is achieved on scientific texts, while reviews prove the most challenging due to short and stylistically homogeneous texts.

\subsection{Experiment 2: Cross-Domain Testing}

We evaluate each model on test sets from all other domains, using thresholds calibrated on the source domain.

\textbf{Results:} Cross-domain performance drops substantially. For example, the model trained on news achieves an F1-score of only 0.69 when tested on scientific texts, compared to 0.85 in-domain. This confirms that perplexity thresholds are domain-dependent and cannot be directly transferred.

\subsection{Experiment 3: Generalization on Mixed-Domain Data}

We combine all domains into a single balanced dataset (7,056 texts total, split 70/15/15 for train/validation/test). A single KenLM model is trained on the combined training set.

\textbf{Results:} The unified model achieves a balanced accuracy of 0.68 and F1-score of 0.68. While this outperforms cross-domain transfer, it still lags behind in-domain performance on individual domains, indicating that domain heterogeneity remains a challenge.

\subsection{Experiment 4: Tokenization and Pruning Optimization}

We experiment with SentencePiece tokenization (vocab\_size=8000) and n-gram pruning ($-prune$ 2 2 2) to improve model stability across domains.

\textbf{Results:} Combined optimization yields a modest improvement of 1-2\% in balanced accuracy (from 0.68 to 0.69), suggesting that while helpful, tokenization and pruning alone are insufficient to fully resolve domain shift.

\subsection{Experiment 5: Ensemble Methods}

We implement four ensemble strategies combining the seven domain-specific models:

\begin{table}[H]
\centering
\caption{Comparison of ensemble methods on the mixed-domain test set.}
\begin{tabular}{|l|c|c|c|c|}
\hline
\textbf{Method} & \textbf{Acc} & \textbf{F1} & \textbf{Balanced Acc} & \textbf{Human/AI} \\
\hline
Simple majority voting & 0.63 & 0.65 & 0.61 & 32\% / 91\% \\
Weighted voting (by val accuracy) & 0.68 & 0.65 & 0.68 & 78\% / 58\% \\
Weighted voting (by balanced acc) & 0.68 & 0.68 & 0.69 & 73\% / 66\% \\
Perplexity-weighted voting & 0.63 & 0.47 & 0.63 & 93\% / 33\% \\
\hline
\end{tabular}
\label{tab:ensemble}
\end{table}

\noindent
\textbf{Observations:}
\begin{itemize}
    \item Simple majority voting strongly favors machine detection (91\% AI recall) but fails on human texts (32\%).
    \item Weighted voting by balanced accuracy achieves the most balanced performance, with human and AI detection rates of 73\% and 66\%, respectively.
    \item Perplexity-weighted voting overfits to human texts, achieving 93\% human recall at the cost of only 33\% AI recall.
\end{itemize}

\subsection{Summary of Key Findings}

\begin{enumerate}
    \item \textbf{Domain shift is significant:} Optimal thresholds vary by a factor of 2-3 across domains (from 8.28 to 12.10).
    \item \textbf{Single-domain models do not generalize:} Cross-domain F1-scores drop by 0.15–0.30 compared to in-domain.
    \item \textbf{Tokenization and pruning provide marginal gains:} Only 1–2\% improvement on mixed-domain data.
    \item \textbf{Weighted ensemble with balanced accuracy is most robust:} Achieves the highest balanced accuracy (0.69) and most balanced human/AI detection rates (73\% / 66\%).
\end{enumerate}


\section{Discussion}\label{sec:disc}

TODO

\section{Conclusion}\label{sec:concl}

TODO


%%%%%%%%%%%%%%%%%%%%%%%%%%%%%%%%%%%%%%%%%%%%%%%%%%%%%%%%%%%%

\bibliographystyle{plain}
\bibliography{references}

%%%%%%%%%%%%%%%%%%%%%%%%%%%%%%%%%%%%%%%%%%%%%%%%%%%%%%%%%%%%

\newpage
\appendix
\section{Appendix / supplemental material}\label{app}

\subsection{Additional experiments / Proofs of Theorems}\label{app:exp}

TODO

\end{document}\end{abstract}

\textbf{Keywords:} Calibrated perplexity
Cross-domain detection
Robustness
Machine-generated text
Text classification
Multilingual detection
Genre adaptation
Statistical methods
Language modeling

\section{Introduction}\label{sec:intro}

Fast improvement and widespread accessibility of Large Language Models (LLMs) have destroyed the lines between human-written and machine-generated text. While this technology offers benefits, it also presents significant disadvantages, including academic dishonesty, disinformation, and the fall of trust in digital content \cite{xu2024detecting, hans2024spotting}. This has created an inevitable need for robust methods to detect AI-generated content.

Current detectors often rely on fine-tuned classifiers or statistical signatures taken from the text \cite{gehrmann2019gltr}. Useful and computationally efficient feature is perplexity, a measure of how "surprising" a text is to a language model. The assumption is that LLMs, which are trained to maximize likelihood, will generate text with lower perplexity compared to human writing. As it is shown in \cite{bakhteev2022automatic}, a significant change comes from distribution shift. The perplexity of a text is not an absolute value; it is highly dependent on the domain (e.g., news, scientific articles, social media) and language it was trained on. A model fine-tuned for one genre often fails on another, causing unreliable detection in real world, where the target domain is unknown.

Recent works have attempted to solve this by developing zero-shot detectors that are more robust to unseen prompts or models. For instance, Binoculars \cite{hans2024spotting} uses a ratio of scores from two different LLMs for calibration. Similarly, research on multilingual and cross-domain detection highlights the necessity of adapting methods to diverse data distributions \cite{gritsay2023artificially}. Despite these advances, a method for calibrating perplexity for each domain itself to make it robust across different genres and languages remains still not-explored.

In this paper, we face the challenge of domain shift and try to avoid its effect on detection. We propose a novel method called calibrated perplexity, which normalizes perplexity scores against domain-specific baselines. The final goal is not to find out one universal calibrated perplexity score, but to get a set of parameters to have an opportunity to learn fastly new zero-shot models. Our approach includeses training KenLM n-gram models on text from multiple domains and using these models to compute a relative, rather than absolute, measure of text predictability. The core novelty of our work lies in transforming a domain-dependent score into a stable feature for classification. To check our method, we conduct experiments across multiple domains and languages, comparing in-domain and out-of-domain detection.


\section{Problem statement}

Let $\mathbf{W}$ be our alphabet, tokens consists from characters from that alphabet.

Let $$\mathbb{D} = \Bigl\{\Big[t_j\Big]_{j=1}^n \Hquad|\Hquad t_j \in \mathbf{W}, n \in \mathbb{N}\Bigr\}$$ be the space of our documents.

We have a set of $N$ documents
$$\mathbf{D} = \bigcup_{i=1}^{N}D^i, D^i \in \mathbb{D}$$

%Each document  $D^i$ consists of $d_i$ tokens $t_1^i, t_2^i, ..., t_{d_i}^i$, where $t_j^i \in \mathbf{W}$ .
We have $K$ classes of our text-generative models, that generated fragments of text in $\mathbf{D}$.

We have $K + 1$ labels, that we use for multiclass classification, where $0$ is the label of human-written fragment, $\{1...K\}$ are the  labels that represent corresponding language model, let $\mathbf{C}$ be our set of labels.

Our model $\mathbf{\phi}$ is a superposition of two mappings, $\mathbf{f}$ and $\mathbf{g}$. Mapping $\mathbf{f}$ is responsible for text segmentation. Mapping $\mathbf{g}$ is responsible for classifying obtained fragments.

$$\mathbf{\phi} : \mathbf{g} \circ \mathbf{f},$$


$$\mathbf{\phi}: \mathbb{D} \rightarrow \mathbb{T}, \qquad \mathbb{T} = \Bigl\{\Big[t_{s_j}, t_{f_j}, C_j\Big]_{j = 1}^{J} \Hquad|\Hquad t_{s_j} = t_{f_{j - 1}},\Hquad s_j \in \mathbb{N}_0,\Hquad f_j \in \mathbb{N}, \Hquad C_j \in \mathbf{C}\Bigr\},$$

where $J$ is a number of fragments in segmentation of the document, $t_{s_j}$ is starting index of the $j$-th fragment,  $t_{f_j}$ is finishing index of the $j$-th fragment,  $C_{j}$ is class of the $j$-th fragment.


To estimate the statistical properties of text fragments, we use a character-aware n-gram language model, KenLM \cite{chen1996empirical}. Let a fragment of text be represented as a sequence of tokens $[t_a, t_{a+1}, ..., t_b]$, where $t_i \in \mathbf{W}$. KenLM provides a function $\mathcal{K}$ that computes the log-probability of this sequence under an n-gram model trained on a specific corpus.

Formally, we define a language model $\mathcal{M}_{\mathcal{T}}$ trained on a text corpus $\mathcal{T} \subset \mathbb{D}$. For a given text fragment $F = [t_{a}, ..., t_{b}]$, the model estimates its probability as the product of conditional probabilities of each token given its $n-1$ predecessors:
\begin{equation}
    P_{\mathcal{M}_{\mathcal{T}}}(F) = \prod_{i=a}^{b} P_{\mathcal{M}_{\mathcal{T}}}(t_i \mid t_{i-1}, ..., t_{i-n+1}).
\end{equation}
KenLM efficiently implements this estimation using modified Kneser-Ney smoothing, which is particularly effective for handling sparse n-gram occurrences.

The **perplexity** of a fragment $F$ is then defined as the inverse geometric mean of the token probabilities:
\begin{equation}
    \text{PPL}_{\mathcal{M}_{\mathcal{T}}}(F) = P_{\mathcal{M}_{\mathcal{T}}}(F)^{-\frac{1}{|F|}} = 2^{-\frac{1}{|F|}\log_2 P_{\mathcal{M}_{\mathcal{T}}}(F)},
\end{equation}
where $|F| = b-a+1$ is the length of the fragment.

This metric quantifies how "surprising" the fragment is to the model $\mathcal{M}_{\mathcal{T}}$; lower perplexity indicates that the fragment is more predictable based on the training corpus $\mathcal{T}$. In our work, we train multiple KenLM models $\{\mathcal{M}_{\mathcal{T}_d}\}$ on different domains $\mathcal{T}_d$ (e.g., news, science, social media). For a given fragment $F$, we compute a set of raw perplexities $\{\text{PPL}_{\mathcal{M}_{\mathcal{T}_d}}(F)\}$. The subsequent mapping $\mathbf{g}$ then uses these values, along with the segmentation information from $\mathbf{f}$, to perform the final classification.

\textbf{Contributions.} Our contributions can be summarized as follows:
\begin{itemize}
    \item We present...
    \item We demonstrate the validity of our theoretical results through empirical studies...
    \item We highlight the implications of our findings for...
\end{itemize}

\textbf{Outline.} The rest of the paper is organized as follows...

\section{Related Work}\label{sec:rw}

\textbf{Topic \#1.}
TODO

\textbf{Topic \#2.}
TODO

\section{Preliminaries}\label{sec:prelim}

\subsection{General notation}

In this section, we introduce the general notation used in the rest of the paper and the basic assumptions. 

\subsection{Assumptions} 

TODO

\section{Method}\label{sec:method}

\section{Experiments}\label{sec:exp}

\subsubsection{Experiment Planning}

\paragraph{Goal of the Computational Experiment}

The goal of the computational experiment is to evaluate the effectiveness of \textit{perplexity-based detection} of AI-generated text using a KenLM language model trained on a balanced dataset of human and machine-generated texts. We try to assess the detection performance on an in-domain test set and to compare different threshold selection strategies (accuracy, F1-score, and ROC-AUC) for the binary classifier.

\paragraph{Data Description}

We use the \textbf{AINL dataset} (\texttt{train-kenlm.csv} and \texttt{dev\_full.csv}), an open-access collection of Russian-language texts with annotations indicating whether a text is human-written or generated by a language model.

\textbf{Data statistics:}

\begin{table}[H]
\centering
\begin{tabular}{|l|l|c|c|c|}
\hline
\textbf{Split} & \textbf{Source} & \textbf{Human Texts} & \textbf{Machine Texts (llama-3.3-70b)} & \textbf{Total} \\
\hline
Train & train-kenlm.csv & 8,769 & 8,700 & 35,158 \\
Validation & dev\_full.csv & $\sim$1,100 & $\sim$1,100 & $\sim$2,200 \\
Test & dev\_full.csv & $\sim$1,100 & $\sim$1,100 & $\sim$2,200 \\
\hline
\end{tabular}
\caption{Data statistics by split.}
\end{table}

\textbf{Data description:}
\begin{itemize}
    \item Texts are in Russian language
    \item Human texts are labeled as `abstract'
    \item Machine texts are labeled with the model name (`llama-3.3-70b')
    \item All texts are preprocessed: filtering empty and short texts ($<100$ characters)
\end{itemize}

\paragraph{Algorithm Configuration}

The base algorithm is \textbf{KenLM}, an n-gram language model with Kneser-Ney smoothing. We train a trigram model ($\text{order}=3$) on the training data. For classification, we use a simple threshold-based detector: a text is classified as machine-generated if its perplexity is below a threshold.

\textbf{Key parameters:}
\begin{itemize}
    \item Model order: 3
    \item Smoothing: Kneser-Ney
    \item Threshold selection: optimal value on validation set
\end{itemize}

\paragraph{Experimental Plan}

\begin{enumerate}
    \item \textbf{Data splitting:}
    \begin{itemize}
        \item Train: all human + all llama-3.3-70b texts from \texttt{train-kenlm.csv}
        \item Validation: 50\% of human + 50\% of llama-3.3-70b from \texttt{dev\_full.csv}
        \item Test: remaining 50\% of human + 50\% of llama-3.3-70b from \texttt{dev\_full.csv}
    \end{itemize}
    
    \item \textbf{Model training:}
    \begin{itemize}
        \item Train KenLM on \texttt{train\_texts.txt} (human + llama texts)
        \item Compute perplexity for all texts
    \end{itemize}
    
    \item \textbf{Threshold selection (on validation set):}
    \begin{itemize}
        \item Find optimal threshold maximizing Accuracy
        \item Find optimal threshold maximizing F1-score
        \item Find optimal threshold by ROC-AUC
    \end{itemize}
    
    \item \textbf{Evaluation (on test set):}
    \begin{itemize}
        \item Apply each threshold to test data
        \item Compute Accuracy, F1-score, ROC-AUC
        \item Compare different threshold strategies
    \end{itemize}
    
    \item \textbf{Analysis:}
    \begin{itemize}
        \item Compare validation vs test performance
        \item Analyze distribution of perplexity for human and machine texts
    \end{itemize}
\end{enumerate}

\paragraph{Expected Tables and Figures}

\textbf{Short list:}
\begin{enumerate}
    \item Table: Optimal thresholds per metric (Accuracy, F1, ROC-AUC)
    \item Figure: Distribution of perplexity (human vs machine) on test set
\end{enumerate}

\textbf{Long list:}

\begin{table}[H]
\centering
\begin{tabular}{|l|l|}
\hline
\textbf{Figure/Table} & \textbf{Description} \\
\hline
Fig. 1 & Accuracy vs threshold curve (X: threshold, Y: accuracy) \\
Fig. 2 & F1-score vs threshold curve (X: threshold, Y: F1) \\
Fig. 3 & Distribution of perplexity with F1-optimal threshold (histogram, human/machine) \\
Fig. 4 & Test set perplexity distributions (histogram, density) \\
Fig. 5 & Test set boxplot (human vs machine) \\
Fig. 6 & ROC curve on test set (X: FPR, Y: TPR) \\
Fig. 7 & Confusion matrix (heatmap) \\
Fig. 8 & Bar chart: validation vs test metrics comparison (F1, Accuracy, ROC-AUC) \\
Fig. 9 & Bar chart: prediction distribution by class \\
Table 1 & Data statistics (splits, counts) \\
Table 2 & Optimal thresholds on validation (F1, Accuracy, Youden) \\
Table 3 & Test results with F1-optimized threshold (Accuracy, F1, ROC-AUC) \\
Table 4 & Comparison: validation vs test metrics (F1, Accuracy, ROC-AUC) \\
\hline
\end{tabular}
\caption{Expected tables and figures.}
\end{table}

\subsubsection{R: Preliminary Report}

\paragraph{Perplexity Statistics}

The perplexity values for validation data show clear separation between classes:

\begin{table}[H]
\centering
\begin{tabular}{|l|c|c|c|}
\hline
\textbf{Class} & \textbf{Mean Perplexity} \\
\hline
Human & 10.97 \\
Machine (llama-3.3-70b) & 5.98 \\
\hline
\end{tabular}
\caption{Perplexity statistics on validation set.}
\end{table}

\paragraph{Threshold Selection}

The optimal threshold maximizing F1-score on validation is:

\[
\text{threshold}_{\text{F1}} = 7.94
\]

At this threshold, validation metrics are:
\begin{itemize}
    \item Accuracy: 0.8776
    \item F1-score: 0.8811
    \item ROC-AUC: 0.9298
\end{itemize}

\paragraph{Test Results}

Applying the F1-optimized threshold to the test set yields:

\begin{table}[H]
\centering
\begin{tabular}{|l|c|}
\hline
\textbf{Metric} & \textbf{Value} \\
\hline
Accuracy & 0.8640 \\
F1-score & 0.8685 \\
ROC-AUC & 0.9306 \\
\hline
\end{tabular}
\caption{Test set results with F1-optimized threshold.}
\end{table}

\paragraph{Validation vs Test Comparison}

\begin{table}[H]
\centering
\begin{tabular}{|l|c|c|c|}
\hline
\textbf{Metric} & \textbf{Validation} & \textbf{Test} & \textbf{Difference} \\
\hline
F1-score & 0.8811 & 0.8685 & -0.0126 \\
Accuracy & 0.8776 & 0.8640 & -0.0136 \\
ROC-AUC & 0.9298 & 0.9306 & +0.0008 \\
\hline
\end{tabular}
\caption{Comparison of validation and test metrics.}
\end{table}

The small difference between validation and test metrics confirms the stability of the model and the representativeness of the data split.

\begin{table}[H]
\centering
\begin{tabular}{|l|c|c|}
\hline
 & \textbf{Predicted Human} & \textbf{Predicted Machine} \\
\hline
\textbf{Actual Human} & 930 & 167 \\
\textbf{Actual Machine} & 135 & 967 \\
\hline
\end{tabular}
\caption{Confusion matrix for test set (F1-optimized threshold).}
\end{table}

\begin{figure}[htbp]
    \centering
    \includegraphics[width=0.7\textwidth]{paper_pic1.png}
\end{figure}

\begin{figure}[htbp]
    \centering
    \includegraphics[width=0.7\textwidth]{paper_pic2.png}
\end{figure}

\paragraph{Interpretation}

The experimental results demonstrate that:
\begin{enumerate}
    \item Perplexity effectively separates human-written and machine-generated texts (validation AUC = 0.93).
    \item The F1-optimized threshold (79.4) provides balanced classification performance.
    \item Test set metrics closely match validation metrics (F1 difference = -0.0126), confirming model stability.
\end{enumerate}

\subsubsection{B: Basic Code}

The basic algorithm implementation is available in the notebook \texttt{exp1.ipynb} in the repository. The key steps include:

\begin{enumerate}
    \item KenLM installation and compilation
    \item Data loading from \texttt{train-kenlm.csv} and \texttt{dev\_full.csv}
    \item Perplexity calculation for each text
    \item Threshold optimization on validation set (by F1-score)
    \item Evaluation on test set
\end{enumerate}

\textbf{Key code sections:}
\begin{itemize}
    \item \texttt{get\_perplexity()}: computes perplexity for a single text
    \item Threshold search loops over possible values
    \item Metrics calculation using \texttt{sklearn}
\end{itemize}

The code is modular and reproducible. The basic model (trigram KenLM with Kneser-Ney smoothing) runs efficiently on CPU, processing approximately 1000 texts per minute.

\vspace{0.5cm}
\noindent\textbf{Repository structure:}
\begin{verbatim}
experiments/
    exp1.ipynb          # main notebook with experiments
    train-kenlm.csv     # training data
    dev_full.csv        # validation and test data
\end{verbatim}


\section{Discussion}\label{sec:disc}

TODO

\section{Conclusion}\label{sec:concl}

TODO


%%%%%%%%%%%%%%%%%%%%%%%%%%%%%%%%%%%%%%%%%%%%%%%%%%%%%%%%%%%%

\bibliographystyle{unsrtnat}
\bibliography{references}

%%%%%%%%%%%%%%%%%%%%%%%%%%%%%%%%%%%%%%%%%%%%%%%%%%%%%%%%%%%%

\newpage
\appendix
\section{Appendix / supplemental material}\label{app}

\subsection{Additional experiments / Proofs of Theorems}\label{app:exp}

TODO

\end{document}\end{abstract}

\textbf{Keywords:} Calibrated perplexity
Cross-domain detection
Robustness
Machine-generated text
Text classification
Multilingual detection
Genre adaptation
Statistical methods
Language modeling

\section{Introduction}\label{sec:intro}

Fast improvement and widespread accessibility of Large Language Models (LLMs) have destroyed the lines between human-written and machine-generated text. While this technology offers benefits, it also presents significant disadvantages, including academic dishonesty, disinformation, and the fall of trust in digital content \cite{xu2024detecting, hans2024spotting}. This has created an inevitable need for robust methods to detect AI-generated content.

Current detectors often rely on fine-tuned classifiers or statistical signatures taken from the text \cite{gehrmann2019gltr}. Useful and computationally efficient feature is perplexity, a measure of how "surprising" a text is to a language model. The assumption is that LLMs, which are trained to maximize likelihood, will generate text with lower perplexity compared to human writing. As it is shown in \cite{bakhteev2022automatic}, a significant change comes from distribution shift. The perplexity of a text is not an absolute value; it is highly dependent on the domain (e.g., news, scientific articles, social media) and language it was trained on. A model fine-tuned for one genre often fails on another, causing unreliable detection in real world, where the target domain is unknown.

Recent works have attempted to solve this by developing zero-shot detectors that are more robust to unseen prompts or models. For instance, Binoculars \cite{hans2024spotting} uses a ratio of scores from two different LLMs for calibration. Similarly, research on multilingual and cross-domain detection highlights the necessity of adapting methods to diverse data distributions \cite{gritsay2023artificially}. Despite these advances, a method for calibrating perplexity for each domain itself to make it robust across different genres and languages remains still not-explored.

In this paper, we face the challenge of domain shift and try to avoid its effect on detection. We propose a novel method called calibrated perplexity, which normalizes perplexity scores against domain-specific baselines. The final goal is not to find out one universal calibrated perplexity score, but to get a set of parameters to have an opportunity to learn fastly new zero-shot models. Our approach includeses training KenLM n-gram models on text from multiple domains and using these models to compute a relative, rather than absolute, measure of text predictability. The core novelty of our work lies in transforming a domain-dependent score into a stable feature for classification. To check our method, we conduct experiments across multiple domains and languages, comparing in-domain and out-of-domain detection.


\section{Problem statement}

Let $\mathbf{W}$ be our alphabet, tokens consists from characters from that alphabet.

Let $$\mathbb{D} = \Bigl\{\Big[t_j\Big]_{j=1}^n \Hquad|\Hquad t_j \in \mathbf{W}, n \in \mathbb{N}\Bigr\}$$ be the space of our documents.

We have a set of $N$ documents
$$\mathbf{D} = \bigcup_{i=1}^{N}D^i, D^i \in \mathbb{D}$$

%Each document  $D^i$ consists of $d_i$ tokens $t_1^i, t_2^i, ..., t_{d_i}^i$, where $t_j^i \in \mathbf{W}$ .
We have $K$ classes of our text-generative models, that generated fragments of text in $\mathbf{D}$.

We have $K + 1$ labels, that we use for multiclass classification, where $0$ is the label of human-written fragment, $\{1...K\}$ are the  labels that represent corresponding language model, let $\mathbf{C}$ be our set of labels.

Our model $\mathbf{\phi}$ is a superposition of two mappings, $\mathbf{f}$ and $\mathbf{g}$. Mapping $\mathbf{f}$ is responsible for text segmentation. Mapping $\mathbf{g}$ is responsible for classifying obtained fragments.

$$\mathbf{\phi} : \mathbf{g} \circ \mathbf{f},$$


$$\mathbf{\phi}: \mathbb{D} \rightarrow \mathbb{T}, \qquad \mathbb{T} = \Bigl\{\Big[t_{s_j}, t_{f_j}, C_j\Big]_{j = 1}^{J} \Hquad|\Hquad t_{s_j} = t_{f_{j - 1}},\Hquad s_j \in \mathbb{N}_0,\Hquad f_j \in \mathbb{N}, \Hquad C_j \in \mathbf{C}\Bigr\},$$

where $J$ is a number of fragments in segmentation of the document, $t_{s_j}$ is starting index of the $j$-th fragment,  $t_{f_j}$ is finishing index of the $j$-th fragment,  $C_{j}$ is class of the $j$-th fragment.


To estimate the statistical properties of text fragments, we use a character-aware n-gram language model, KenLM \cite{chen1996empirical}. Let a fragment of text be represented as a sequence of tokens $[t_a, t_{a+1}, ..., t_b]$, where $t_i \in \mathbf{W}$. KenLM provides a function $\mathcal{K}$ that computes the log-probability of this sequence under an n-gram model trained on a specific corpus.

Formally, we define a language model $\mathcal{M}_{\mathcal{T}}$ trained on a text corpus $\mathcal{T} \subset \mathbb{D}$. For a given text fragment $F = [t_{a}, ..., t_{b}]$, the model estimates its probability as the product of conditional probabilities of each token given its $n-1$ predecessors:
\begin{equation}
    P_{\mathcal{M}_{\mathcal{T}}}(F) = \prod_{i=a}^{b} P_{\mathcal{M}_{\mathcal{T}}}(t_i \mid t_{i-1}, ..., t_{i-n+1}).
\end{equation}
KenLM efficiently implements this estimation using modified Kneser-Ney smoothing, which is particularly effective for handling sparse n-gram occurrences.

The **perplexity** of a fragment $F$ is then defined as the inverse geometric mean of the token probabilities:
\begin{equation}
    \text{PPL}_{\mathcal{M}_{\mathcal{T}}}(F) = P_{\mathcal{M}_{\mathcal{T}}}(F)^{-\frac{1}{|F|}} = 2^{-\frac{1}{|F|}\log_2 P_{\mathcal{M}_{\mathcal{T}}}(F)},
\end{equation}
where $|F| = b-a+1$ is the length of the fragment.

This metric quantifies how "surprising" the fragment is to the model $\mathcal{M}_{\mathcal{T}}$; lower perplexity indicates that the fragment is more predictable based on the training corpus $\mathcal{T}$. In our work, we train multiple KenLM models $\{\mathcal{M}_{\mathcal{T}_d}\}$ on different domains $\mathcal{T}_d$ (e.g., news, science, social media). For a given fragment $F$, we compute a set of raw perplexities $\{\text{PPL}_{\mathcal{M}_{\mathcal{T}_d}}(F)\}$. The subsequent mapping $\mathbf{g}$ then uses these values, along with the segmentation information from $\mathbf{f}$, to perform the final classification.

\textbf{Contributions.} Our contributions can be summarized as follows:
\begin{itemize}
    \item We present...
    \item We demonstrate the validity of our theoretical results through empirical studies...
    \item We highlight the implications of our findings for...
\end{itemize}

\textbf{Outline.} The rest of the paper is organized as follows...

\section{Related Work}\label{sec:rw}

\textbf{Topic \#1.}
TODO

\textbf{Topic \#2.}
TODO

\section{Preliminaries}\label{sec:prelim}

\subsection{General notation}

In this section, we introduce the general notation used in the rest of the paper and the basic assumptions. 

\subsection{Assumptions} 

TODO

\section{Method}\label{sec:method}

\section{Experiments}\label{sec:exp}

To verify the theoretical estimates obtained, we conducted a detailed empirical study...

\section{Discussion}\label{sec:disc}

TODO

\section{Conclusion}\label{sec:concl}

TODO


%%%%%%%%%%%%%%%%%%%%%%%%%%%%%%%%%%%%%%%%%%%%%%%%%%%%%%%%%%%%

\bibliographystyle{unsrtnat}
\bibliography{references}

%%%%%%%%%%%%%%%%%%%%%%%%%%%%%%%%%%%%%%%%%%%%%%%%%%%%%%%%%%%%

\newpage
\appendix
\section{Appendix / supplemental material}\label{app}

\subsection{Additional experiments / Proofs of Theorems}\label{app:exp}

TODO

\end{document}
